\documentclass[11pt,a4paper]{article}

% Packages
\usepackage[utf8]{inputenc}
\usepackage[T1]{fontenc}
\usepackage{amsmath}
\usepackage{amssymb}
\usepackage{amsthm}
\usepackage{enumitem}
\usepackage{geometry}
\usepackage{graphicx}
\usepackage{hyperref}
\usepackage{listings}
\usepackage{xcolor}

% Page geometry
\geometry{margin=1in}

% Code listing settings
\lstset{
  basicstyle=\ttfamily\small,
  breaklines=true,
  frame=single,
  backgroundcolor=\color{gray!10}
}

% Custom commands
\newcommand{\code}[1]{\texttt{#1}}

% Document metadata
\title{Inflation Simulation in ORE}
\date{2026-02-09}

\begin{document}

\maketitle

\begin{abstract}
This document describes the desired theory and practical approach for simulating inflation in ORE with the CrossAssetModel. It focuses on how seasonality, day-count differences and observation lags should be handled.
\end{abstract}

\section{Conventions and Terminology}

Before discussing simulation, we establish the t=0 market curve conventions that must be preserved throughout the exposure simulation.

\subsection{Key Dates and Lags}

\begin{itemize}[leftmargin=*]

\item \textbf{Reference date:} The valuation date, denoted \code{today} or $t=0$.

\item \textbf{Curve base date ($\text{baseDate}_{\text{curve}}$):} The date of the last known published CPI fixing. Always the \textbf{first day of an inflation period} (typically the 1st of a month for monthly frequency indices). This is the anchor for the zero-inflation curve.

\item \textbf{Simulation lag (simLag):} The difference between the curve base date and reference date.
  
  Example: today = 9-Feb-2026 and last published CPI is Dec-2025, so simLag $\approx$ 2M + 8d.

\item \textbf{Published fixings:} CPI indices publish one fixing per calendar month, always assigned to the \textbf{first day of the month} (the start of the inflation period). There are no intra-month fixings. Forward CPI are flat throughout the month.

\item \textbf{Flat forward rule:} Between published fixings (within a calendar month), forward CPI values remain \textbf{constant} -- the forward CPI for any date within month M equals the forward CPI for the 1st of month M.

\item \textbf{Interpolation:} When instruments require interpolated CPI (e.g., \code{CPI::Linear} interpolation), the interpolation is performed \textbf{by the instrument/index} between the two monthly fixings that bracket the observation date. The interpolation uses the current calendar date within the month (not lagged time).
\end{itemize}

\subsection{T0 Zero-Inflation Curves}

Zero-coupon inflation index swaps (ZCIIS) are used for bootstrapping the zero-inflation curve:

\begin{itemize}[leftmargin=*]
\item \textbf{Observation lag (obsLag):} A \textbf{contractual convention} specified in zero-coupon inflation swap (ZCIIS) instruments. For a swap with payment at date T, the CPI observation occurs at \code{T - obsLag}, then moved to the start of the inflation period. Example: obsLag = 3 months.
\end{itemize}

Assume:
\begin{itemize}
\item \textbf{Today} (reference date): 9-Feb-2026
\item \textbf{ZCIIS tenor:} 2 years
\item \textbf{Observation lag:} 3 months
\item \textbf{Index frequency:} monthly (first-of-month fixings)
\end{itemize}

Then:

\begin{itemize}[leftmargin=*]
\item \textbf{ZCIIS base date ($t_\text{base} = \text{baseDate}_{\text{ZCIIS}}$):}
    \begin{itemize}
    \item First, lag today by obsLag: 9-Feb-2026 - 3M = 9-Nov-2025
    \item Then move to first-of-month: \textbf{1-Nov-2025}
    \end{itemize}

\item \textbf{Maturity (payment date T):} 9-Feb-2026 + 2Y = 9-Feb-2028

\item \textbf{Observation date ($t_\text{obs}$):}
    \begin{itemize}
    \item First, lag maturity by obsLag: 9-Feb-2028 - 3M = 9-Nov-2027
    \item Then discretized to first-of-month: \textbf{1-Nov-2027}
    \end{itemize}

\item \textbf{ZCIIS structure:} A ZCIIS with maturity T and start date \code{today} pays $I(T_{\text{obs}}) / I(\text{baseDate}_{\text{ZCIIS}}) - 1$ at maturity T, where $T_{\text{obs}} = \text{firstOfInflationPeriod}(T - \text{obsLag})$ and $\text{baseDate}_{\text{ZCIIS}} = \text{firstOfInflationPeriod}(\text{today} - \text{obsLag})$.

\item \textbf{Pillar dates:} For a ZCIIS with tenor (e.g., 1Y, 2Y, 5Y), the pillar observation date is computed as:
\begin{lstlisting}
maturity = today + tenor
t_obs = firstOfInflationPeriod(maturity - obsLag)
\end{lstlisting}
The curve is bootstrapped with pillars at these \code{t\_obs} dates.
\end{itemize}

\subsection{Simplified Timeline}

Below is a conceptual timeline illustrating the key dates and lags for a ZCIIS with 2Y tenor and 3M observation lag. The example uses the dates from sections 1.1 and 1.2 (today = 9-Feb-2026, last CPI = Dec-2025), where the simulation lag and observation lag differ:

\begin{verbatim}
Time axis (calendar dates):

  |------|------|------|---------- ... ----------|------|------|------|
  Nov    Dec    Jan    Feb                       Nov    Dec    Jan    Feb
  2025   2025   2026   2026                      2027   2027   2028   2028
  ^      ^             ^                         ^                    ^
baseDate baseDate      today                     t_obs              maturity (T)
_ZCIS   _curve         (reference                (ZCIIS              (payment date)
         (curve         date)                     observation)
          anchor)
1-Nov   1-Dec          9-Feb                     1-Nov              9-Feb
-2025   -2025          -2026                     -2027              -2028

  <-------obsLag------->                          <-------obsLag------->
           = 3M                                            = 3M
         <--simLag----->
           ~ 2M + 8d
\end{verbatim}

\noindent\textbf{Curve anchor:}  
$\text{baseDate}_\text{curve} = \text{1-Dec-2025}$ (last known published CPI fixing)

\noindent\textbf{ZCIIS base:}  
$\text{baseDate}_\text{ZCIIS} = \text{1-Nov-2025}$ (firstOfMonth(today - obsLag))

\noindent\textbf{Pillar date:}  
$t_\text{obs} = \text{1-Nov-2027}$ (firstOfMonth(maturity - obsLag))

\noindent\textbf{Fair ZCIIS rate:}
\begin{equation}
\frac{I(t_\text{obs})}{I(\text{baseDate}_\text{ZCIIS})} = (1 + z_\text{fair})^{\tau_0} \quad \text{paid at maturity } T
\end{equation}
where $\tau_0 = \text{yearFraction}(\text{baseDate}_\text{ZCIIS}, t_\text{obs})$

\noindent\textbf{Curve zero rate:} given by this equation:
\begin{equation}
I(t_\text{obs}) = I(\text{baseDate}_\text{curve}) \cdot (1 + z_0)^{\tau_1}
\end{equation}
where $\tau_1 = \text{yearFraction}(\text{baseDate}_\text{curve}, t_\text{obs})$

\noindent\textbf{Key observations:}
\begin{itemize}
\item \code{simLag = 2M + 8d} (today is about 2 months and 8 days after the curve anchor $\text{baseDate}_{\text{curve}}$ = 1-Dec-2025).
\item \code{obsLag = 3M} (contractual lag in the ZCIIS). \textbf{simLag $\neq$ obsLag} in general.
\item The $\text{baseDate}_{\text{curve}}$ (1-Dec-2025, last known fixing) differs from $\text{baseDate}_{\text{ZCIIS}}$ (1-Nov-2025, derived from obsLag). The curve stores zero rates relative to $\text{baseDate}_{\text{curve}}$, while swap payoffs reference $\text{baseDate}_{\text{ZCIIS}}$.
\item The ZCIIS pillar is at $T_{\text{obs}}$ = 1-Nov-2027.
\item During any calendar month, all forward CPIs for dates within that month are identical.
\item Interpolation methods: linear in zeroRate or log-linear in CPI
\end{itemize}

\subsection{Multiplicative Price Seasonality}

Consumer price indices exhibit deterministic seasonal patterns. The observed (market) CPI is modelled as the product of a \textbf{deseasonalised index} $I(d)$ and a \textbf{seasonal factor} $s(d)$:

\begin{equation}
\text{CPI}_{\text{market}}(d) = s(d) \cdot I(d)
\end{equation}

The seasonal factors $\{s_0, s_1, \ldots, s_{11}\}$ are periodic (one per calendar month, repeating each year) and purely deterministic. They capture the ratio between the observed price level and the underlying trend, so that $I(d)$ is free of calendar-driven oscillations.

\subsubsection{The Seasonality-Corrected Zero Rate}

The zero-inflation curve stores deseasonalised zero rates $z_0$. By definition, the deseasonalised index grows from $\text{baseDate}_{\text{curve}}$ to observation date $d$ as:

\begin{equation}
I(d) = I(\text{baseDate}_{\text{curve}}) \cdot (1 + z_0)^{\tau}, \qquad \tau = \text{infDc.yearFraction}(\text{baseDate}_{\text{curve}}, d)
\end{equation}

Multiplying both sides by the respective seasonal factors gives the market CPI relationship:

\begin{equation}
s(d) \cdot I(d) = s(\text{baseDate}_{\text{curve}}) \cdot I(\text{baseDate}_{\text{curve}}) \cdot (1 + z_0)^{\tau}
\end{equation}

We want the market-observed (seasonally-adjusted) zero rate $z$ to satisfy:

\begin{equation}
s(d) \cdot I(d) = s(\text{baseDate}_{\text{curve}}) \cdot I(\text{baseDate}_{\text{curve}}) \cdot (1 + z)^{\tau}
\end{equation}

This gives us the seasonality adjusted zero rate $z$:

\begin{equation}
\boxed{1 + z = (1 + z_0) \cdot \left(\frac{s(d)}{s(\text{baseDate}_{\text{curve}})}\right)^{1/\tau}}
\end{equation}

\subsubsection{Seasonality and the Simulation Drift}

Because seasonality is \textbf{deterministic}, it does not need to be included in the stochastic simulation. The model simulates only the deseasonalised index $I(t)$; seasonality is applied afterwards -- at the publishing stage -- exactly as it is at $t=0$. This means:

\begin{itemize}
\item The model drift is computed using the \textbf{deseasonalised} zero-inflation curve (i.e. after stripping seasonality). The stochastic dynamics and drift corrections (e.g. the JY real-rate/index cross-terms) operate on smooth, deseasonalised quantities only.
\item After simulating to time $t_S$, the deseasonalised CPI $I(t_S)$ is converted to a market CPI by multiplying with $s(d)/s(\text{baseDate}_{\text{curve}})$ (the seasonal ratio relative to the curve anchor), and deseasonalised zero rates are converted to market zero rates using the formula from section 1.4. This is identical to how the $t=0$ curve applies seasonality when queried.
\end{itemize}

\textbf{Problem: drift computation with a step-function seasonality curve.} In a discrete-time simulation scheme (e.g. Euler), the drift at each time step involves the instantaneous forward rate of the zero-inflation curve. With seasonality attached, the market zero-inflation curve is a \textbf{step function} in time (it jumps at each month boundary due to the multiplicative seasonal factors). Numerical differentiation of a step function produces spikes at month boundaries and zeros elsewhere.

The solution is to compute the drift from the \textbf{deseasonalised} curve, which is smooth and well-behaved under numerical differentiation. Seasonality is then re-applied only at the output/publishing stage. This separation is essential for stable discrete-time simulation.

\subsection{Remarks About the Current Implementation}

\begin{itemize}[leftmargin=*]
\item \code{ZeroInflationCurve} stores deseasonalised zero rates.
\item \code{zeroRate(d)} discretizes $d$ to the first day of the inflation period, applies seasonality, and returns the seasonality-corrected zero rate.
\item \code{zeroRate(t)} returns the continuous, deseasonalised zero rate (interpolated within a month using the same scheme).
\item \code{correctZeroRate(...)} in the seasonality curve also applies discretization.
\item \code{ZeroInflationIndex}:
  \begin{itemize}
  \item Adding a fixing stores a flat fixing for the full inflation period (from period start to period end).
  \item Past fixings are discretized and moved to the first day of the inflation period.
  \item Forecast fixings use the base fixing (discrete), discretize the fixing date, and compute
    \begin{equation}
    \text{baseFixing} \cdot (1 + z(\text{obsDate}))^{\tau}
    \end{equation}
    where $\tau = \text{yearFraction}(\text{baseDate}_\text{curve}, \text{obsDate})$.
  \end{itemize}
\end{itemize}

\section{Inflation Simulation}

Our goal is to produce simulated inflation paths that are consistent with the model state and the t=0 zero-inflation curve, reproducing the correct forward CPI forecasts. The simulated zero-inflation curve must satisfy the \textbf{martingale property}: when expressed in terms of the numeraire, the expected value of any inflation-linked cashflow must equal its value computed from the t=0 market curve.

\subsection{Continuous-Time Evolution and Simulation Lag}

The inflation models (JY, DK) evolve the index in \textbf{continuous time}: the model state is defined for every real $t \geq 0$, not only at the start of an inflation period. At model time $t$, the model implies a deseasonalised CPI level $I(t)$ (e.g. under JY: $I(t) = \exp(y(t))$; under DK the mapping differs but the principle is the same).

The relationship between model time and calendar dates is set up so that model time $t = 0$ corresponds to the curve's $\text{baseDate}_{\text{curve}}$ (not to today). In other words, model time $t_0 = \text{today} - \text{baseDate}_{\text{curve}}$ (in days).

When we simulate to calendar date $S$, the model time is simply:

\begin{equation}
t_S = S - \text{simLag}
\end{equation}

and the model implies the deseasonalised CPI $I(t_S)$ at calendar date $S$ -- in continuous time, \textbf{not discretized} to the inflation periods.

We keep the simulation lag constant throughout the simulation.

\textbf{The firstOfMonth rounding is not part of the model.} It only arises when we convert continuous model output into discrete monthly fixings for pricing. Specifically:
\begin{itemize}
\item Instruments request CPI fixings at first-of-month dates (the inflation period convention).
\item For the simulated zero-inflation curve, pillar dates $T_{\text{obs}}$ must be aligned to first-of-month to ensure \textbf{consistency with the t=0 curve}. For example: if we simulate 1Y forward, the 1Y pillar on the simulated curve must correspond to the same first-of-month observation date as the 2Y pillar on today's curve. If pillar dates were not aligned, interpolation between pillars would introduce spurious differences between the simulated forward CPI and the t=0 forward CPI.
\end{itemize}

\subsection{Martingale Property for Inflation-Indexed Cashflows}

Consider an inflation-indexed cashflow that pays $I(T_{\text{obs}}) / I(\text{baseDate}_{\text{ZCIIS}})$ at maturity $T$. The two date conventions from section 1.2 must be distinguished:

\begin{itemize}[leftmargin=*]
\item \textbf{$\text{baseDate}_{\text{ZCIIS}}$, $T_{\text{obs}}$:} the CPI observation dates of the ZCIIS, derived from the \textbf{contractual obsLag}:
\begin{equation}
T_{\text{obs}} = \text{firstOfInflationPeriod}(T - \text{obsLag}), \quad \text{baseDate}_{\text{ZCIIS}} = \text{firstOfInflationPeriod}(\text{today} - \text{obsLag})
\end{equation}

\item \textbf{$\text{baseDate}_{\text{curve}}$:} the curve's anchor date (last known published CPI fixing), derived from \textbf{simLag}:
\begin{equation}
\text{baseDate}_{\text{curve}} = \text{today} - \text{simLag}
\end{equation}
\end{itemize}

In general $\text{baseDate}_{\text{ZCIIS}} \neq \text{baseDate}_{\text{curve}}$ because $\text{obsLag} \neq \text{simLag}$ (section 1.3).

We require $T_{\text{obs}}$ to lie in the future of the simulated curve at time $S$. The simulated curve's base date moves to $\text{baseDate}_{\text{curve},S} = S - \text{simLag}$ (section 2.1), so the condition is $S - \text{simLag} < T_{\text{obs}}$, i.e. $S < T_{\text{obs}} + \text{simLag}$.

\subsubsection{Inflation Growths (Zero-Inflation Curve Encoding)}

At $t=0$, we compute the forward CPI growth from $\text{baseDate}_{\text{curve}}$ to $T_{\text{obs}}$:

\begin{equation}
\frac{I(T_{\text{obs}})}{I(\text{baseDate}_{\text{curve}})} = (1 + z_0(T_{\text{obs}}))^{\tau_0}, \qquad \tau_0 = \text{infDc.yearFraction}(\text{baseDate}_{\text{curve}}, T_{\text{obs}})
\end{equation}

At simulation time $S$, the simulated curve's base date \textbf{moves} to $\text{baseDate}_{\text{curve},S} = S - \text{simLag}$ (cf. \code{ZeroInflationCurveObserverMoving}). The simulated zero rate $z_S$ is defined relative to this new anchor:

\begin{equation}
\frac{I_S(T_{\text{obs}})}{I_S(\text{baseDate}_{\text{curve},S})} = (1 + z_S(T_{\text{obs}}))^{\tau_S}, \qquad \tau_S = \text{infDc.yearFraction}(\text{baseDate}_{\text{curve},S}, T_{\text{obs}})
\end{equation}

Note $\tau_S < \tau_0$ since $\text{baseDate}_{\text{curve},S} > \text{baseDate}_{\text{curve}}$.

\subsubsection{Martingale Condition}

The bond payoff $I(T_{\text{obs}}) / I(\text{baseDate}_{\text{ZCIIS}})$ involves $\text{baseDate}_{\text{ZCIIS}}$ (the contractual CPI base) which is a \textbf{known constant} (a published fixings), so it cancels between the two sides. Writing the martingale in terms of curve quantities only, we get:

\begin{equation}
\boxed{\mathbb{E}^{\mathbb{Q}}_0 \left[ I_S(\text{baseDate}_{\text{curve},S}) \cdot (1 + z_S(T_{\text{obs}}))^{\tau_S} \cdot \frac{P(S,T)}{N(S)} \right] = I(\text{baseDate}_{\text{curve}}) \cdot (1 + z_0(T_{\text{obs}}))^{\tau_0} \cdot P(0,T)}
\end{equation}

where:
\begin{itemize}
\item $\text{baseDate}_{\text{curve}} = \text{today} - \text{simLag}$: the t=0 curve base date (last known CPI at $t=0$)
\item $\text{baseDate}_{\text{curve},S} = S - \text{simLag}$: the simulated curve's base date at time $S$
\item $z_0(T_{\text{obs}})$: the t=0 zero rate from $\text{baseDate}_{\text{curve}}$ to $T_{\text{obs}}$
\item $z_S(T_{\text{obs}})$: the simulated zero rate from $\text{baseDate}_{\text{curve},S}$ to $T_{\text{obs}}$
\item $\tau_0 = \text{infDc.yearFraction}(\text{baseDate}_{\text{curve}}, T_{\text{obs}})$
\item $\tau_S = \text{infDc.yearFraction}(\text{baseDate}_{\text{curve},S}, T_{\text{obs}})$
\item $I_S(\text{baseDate}_{\text{curve},S})$: the model-implied CPI at the simulated curve's base date
\item $P(S,T)$: the nominal zero-coupon bond price at $S$ for maturity $T$
\item $N(S)$: the nominal numeraire at $S$
\end{itemize}

\subsection{Seasonality in the Simulation}

This follows directly from sections 1.4 and 2.1:

\begin{itemize}
\item The model simulates only the \textbf{deseasonalised} index $I(t)$. The seasonal factors $s(d)$ are deterministic and applied at the publishing stage.
\end{itemize}

\subsection{Scenario Generation: Computing Zero Rates for the Simulated Curves}

Assume there are two model-dependent functions:
\begin{enumerate}
\item $I(S, \text{state}(S))$: computes the deseasonalised CPI at time $t_S = S - \text{simLag}$, given the time $S$ and the model state at time $S$.
\item $G(S, T, \text{state}(S))$: computes the deseasonalised forward growth from model time $S - \text{simLag}$ to $T - \text{simLag}$, given the simulated state at time $S$.
\end{enumerate}

Then the simulated zero rate is:
\begin{equation}
z_S(T_{\text{obs}}) = G(S, T_{\text{obs}} + \text{simLag}, \text{state}(S))^{1/\tau_S} - 1
\end{equation}

and the simulated base CPI is:
\begin{equation}
I_S(S - simLag) = I(S, \text{state}(S))
\end{equation}

\textbf{Note:} The observation date, simulation base date and simLag needs to converted to times with the simulation grid day counter, to be consistent to the generated paths.

\subsection{Remarks:}
\begin{itemize}[leftmargin=*]
\item \textbf{Day counter alignment:} If the simulation grid day counter differs from the inflation day counter (\code{infDc}): (i) compute \code{simLag} using the simulation grid day counter, (ii) convert simulation times to \code{infDc} when querying zero rates from the t=0 curve.

\item \textbf{Pillar dates:} Must be discretized to first-of-month as described in section 2.1 to ensure consistency with the t=0 curve structure.

\item \textbf{Moving base date:} The simulated zero curve requires a continuously moving base date $\text{baseDate}_{\text{curve},S} = S - \text{simLag}$, which generally does \textbf{not} align with start date of the inflation periods. This differs from the t=0 curve where the base date is fixed at the last published first-of-month fixing.

\item \textbf{Fixings at arbitrary dates:} The current \code{ZeroInflationIndex} implementation stores only flat monthly fixings (constant for the entire inflation period, section 1.6). However, for simulation, the curve's base date $\text{baseDate}_{\text{curve},S}$ can fall on any calendar date (e.g., 9-Feb-2026), not just the first of a month. To compute forward CPIs correctly using the simulated zero rates, the inflation index must provide the \textbf{continuous} (non-discretized) CPI value at $\text{baseDate}_{\text{curve},S} = S - \text{simLag}$. Specifically:
  \begin{itemize}
  \item The simulated curve computes zero rates $z_S(T_{\text{obs}})$ relative to the continuous base date $\text{baseDate}_{\text{curve},S}$.
  \item Forward CPI computation uses: $I_S(T_{\text{obs}}) = I_S(\text{baseDate}_{\text{curve},S}) \cdot (1 + z_S(T_{\text{obs}}))^{\tau_S}$
  \item Therefore, $I_S(\text{baseDate}_{\text{curve},S})$ must be stored at the exact arbitrary date, not rounded to the first of the month.
  \end{itemize}

\end{itemize}


\end{document}
